\chapter{Experimentos y Análisis de Resultados}\label{experimentos}

En este capítulo se debe presentar la configuración experimental, las pruebas propuestas y aplicadas, así como el análisis detallado de los resultados alcanzados.

Adicionalmente a lo anterior, se deberá incluir una sección titulada \textbf{Análisis de Madurez Tecnológica}, en la cual se deberá:
\begin{itemize}
	\item Presentar una revisión del nivel TRL (Technology Readiness Level) del resultado/producto alcanzado. 
	\item Analizar porqué el resultado/producto alcanzado está en un nivel u otro y que ruta se espera seguir para continuar evolucionándolo.
	\item Como recomendación, para los proyectos que alcancen una alta madurez tecnológica, se podría relacionar qué ventajas se tiene frente a productos cercanos para establecer de alguna manera si es una innovación, identificando qué barreras se encuentran para entrar al mercado, económicas y legales (normativas, estándares, etc.), entre otras.
\end{itemize}

Como referencia, se puede seguir el documento adjunto sobre \textbf{Níveles de Madurez Tecnológica del Ministerio de Ciencia y Tecnología (MinCiencias)}.





\textcolor{red}{Se recomienda tener una extensión de $10$ páginas.}
